\documentclass[10pt,twocolumn]{article}
\usepackage[margin=0.8in]{geometry}
\usepackage{amsmath,amssymb,booktabs,graphicx,xcolor,hyperref}
\usepackage[numbers]{natbib}
\newcommand{\pending}[1]{\colorbox{yellow!40}{\textsf{\small PENDING: #1}}}

\title{Look Where the Grasp Fails: Calibration-Aware Active Viewpoint Selection\\
over a Multimodal 6D Pose Belief}
\author{Param Madan}
\date{October 2026}

\begin{document}
\maketitle

\begin{abstract}
A robot that may move its camera before grasping can choose views that reduce pose
uncertainty, or views that reduce the chance that the grasp fails. We test, in simulation and
under pre-registered protocols, whether selecting views by \emph{expected grasp failure} under a
temporal, multimodal SE(3) particle belief beats information objectives (hypothesis entropy,
expected rotational covariance) at equal camera budgets. A task-risk objective that treats the camera
calibration as exact does \emph{not} help: it predicts near-zero risk and never moves (null,
150 paired episodes). The limiting error is a per-episode extrinsic calibration bias that is fixed
in the camera frame, so its world-frame effect depends on the view. Modelling that bias in the
predicted grasp loss (D2) raises physical grasp success from 0.52 to 0.67 and beats both
information objectives ($+0.13$, Holm $p=0.0078$). Online self-calibration from the views D2
collects adds a further $+0.13$ (0.64 $\rightarrow$ 0.77, $p=0.0034$). The gains come mostly from
view \emph{diversity} plus bias estimation: once self-calibrated, D2 is statistically
indistinguishable from random views. A final pre-registered study on 299 new paired episodes
replicates both effects and shows the full system beating a Kalman-tracking baseline (0.773 vs
0.488, $+0.28$). With a fixed camera, single-frame, Kalman and belief estimators are
indistinguishable (0.48--0.50), so the gain comes from where the camera looks and from the bias
those views reveal. After per-sensor-configuration scaling on held-out data, the belief's
position credible regions reach 50/80/95\,\% coverage within 5 points. All results are
simulation only.
\end{abstract}

\section{Introduction}
Active perception for manipulation usually picks the next view to reduce uncertainty: the entropy
of pose hypotheses~\cite{liu2025activepose} or the predicted rotational
covariance~\cite{jin2025poseflow}. Uncertainty is a proxy; what matters is whether the grasp
succeeds. The principle that observation should be chosen for its effect on task outcome is
established in belief-space task planning~\cite{curtis2024potamp}. We do not claim it. Our
question is narrower and empirical: \emph{at the level of camera viewpoints over a continuous
multimodal 6D belief, at equal budgets and with paired seeds, does choosing views by expected
grasp loss beat choosing them by information gain, and why?}

\paragraph{Contributions.}
\begin{enumerate}
\item A controlled, pre-registered comparison of task-risk and information-gain view selection at
equal camera budgets, scored by executing grasps on a simulated Franka Panda.
\item A negative result with a mechanism: a task-risk objective that treats calibration as exact
ties with a fixed camera, because the error that limits grasping is view-dependent calibration
bias that none of the objectives model.
\item A calibration-aware task-risk objective (D2) and online self-calibration from the
collected views, each confirmed on held-out seeds and replicated in an eight-arm component ablation.
\item An open, reproducible stack: simulation, perception, belief tracker, policies, grasp
executor, C++/ONNX/TensorRT inference paths, and an experiment registry in which every claim is
traceable to a command, commit and seed range.
\end{enumerate}

\section{Related work}
\textbf{Multimodal 6D pose.} PoseRBPF~\cite{deng2019poserbpf} tracks a Rao-Blackwellised
particle belief that handles symmetry and occlusion; SE(3)-PoseFlow~\cite{jin2025poseflow}
samples pose distributions with flow matching. Our belief tracker is a baseline-level component,
not a contribution over these.
\textbf{Active pose estimation.} ActivePose~\cite{liu2025activepose} selects views by hypothesis
entropy plus a VLM ambiguity score; SE(3)-PoseFlow minimises predicted rotational covariance over
candidate views. Neither links the view objective to grasp outcome. Our C1 and C2 comparators
implement these two objectives (without the VLM term).
\textbf{Task-driven observation.} POTAMP~\cite{curtis2024potamp} shows information gathering
emerging from goal-directed belief-space planning over abstract task beliefs; it does not compare
viewpoint objectives.
\textbf{Pose estimators.} FoundationPose~\cite{wen2024foundationpose} is the strongest published
estimator; we did not run it (Section~\ref{sec:limits}).

\section{System}
\textbf{Simulation.} MuJoCo~\cite{todorov2012mujoco} tabletop with a Franka Panda (Menagerie),
randomised primitive objects and mugs, a fixed overhead and a mobile RGB-D camera. Noise presets
model depth noise, dropout and a per-episode extrinsic calibration error
(${\sim}1^\circ$, 5\,mm). Seeds are split by range (train / calibration / val / test) and
inference code may not import simulator state (enforced by an import-leakage test).

\textbf{Perception.} RANSAC table plane, depth-continuous segmentation with a hue vote, and a
robust SDF fit of the matched primitive with a support prior. A small point network proposes
pose hypotheses that the SDF fit refines (EXP-M3-001: success on views under 50\,\% visible rises
from 0.853 to 0.985).

\textbf{Belief.} Per object, a particle filter on SE(3)$\times\mathbb{R}^3$ with a two-sided
likelihood (surface agreement and free-space violations), progressive
correction~\cite{oudjane2000progressive} and mode-preserving
resampling~\cite{vermaak2003mixture}. Against a Kalman tracker on the same detections it raises
rotation success on partial views from 64.8\,\% to 87.0\,\% ($p=0.0074$) with no MOTA loss
(EXP-M4-001).

\textbf{View policies} (budget: 2 camera moves after one shared warm-up frame; 24 candidate views).
\emph{fixed}; \emph{random}; \emph{C1} mode-weight entropy; \emph{C2} expected rotational variance;
\emph{D} expected grasp failure of the grasp the system would execute after the next view,
estimated from belief particles with the tracker's own observation model; \emph{D2} as D but with
the camera's calibration error as a latent in the predicted loss, so views whose biases cancel
lower the expected failure, plus a travel cost ($\lambda=0.005$).

\textbf{Self-calibration.} A linear-Gaussian MAP estimate of each camera's extrinsic error from
matched track centres across views. Only the view-varying part is identifiable (a gauge freedom);
the grasp centre is corrected by the estimated bias over the views that saw the target.

\textbf{Grasping.} A top-down parallel-jaw grasp planned from the target belief's MAP mode,
executed on the simulated Panda; success means the target object is lifted more than 8\,cm
(evaluator reads simulator state).

\section{Experimental protocol}
Every confirmatory experiment was written into the registry, with frozen commit, seed range,
primary test and analysis command, before it ran. Episodes are paired: each policy replays the
same seed, so initial observations are bit-identical. Primary tests are exact two-sided McNemar
tests on paired outcomes with Holm correction across a family; we report bootstrap 95\,\% CIs.
Development used only val seeds. Regime for all policy studies: mug scenes, harsh noise, mobile
camera only, views at 12/25/45$^\circ$ elevation, physical scoring, one attempt.

\section{Results}
\subsection{Task risk that treats calibration as exact does not move (EXP-M5-001)}
On 150 test episodes D tied with fixed (0.487 vs 0.487) and lost slightly to C1/C2 (0.507;
$\Delta=-0.020$, Holm $p=1.0$): \textbf{H3 not supported}. D predicted near-zero risk under the
belief and stayed put. Unexpectedly, random views won (0.627; $p=0.0019$ vs D). Random's grasp
centre error was lower in 73\,\% of episodes. The extrinsic error is fixed in the camera frame, so
moving the camera re-samples its world-frame effect and diverse views average it out.

\subsection{Calibration-aware task risk (EXP-M5-002)}
D2 was designed on val from this mechanism and re-registered.
\begin{table}[h]\centering\small
\begin{tabular}{lccc}\toprule
Policy & Success & Centre err. & Views\\\midrule
fixed & 0.520 & 18.3\,mm & 1.00\\
D (task risk) & 0.520 & 18.3\,mm & 1.00\\
C1 entropy & 0.533 & 17.1\,mm & 1.06\\
C2 covariance & 0.533 & 17.1\,mm & 1.06\\
random & 0.633 & 14.4\,mm & 1.95\\
\textbf{D2} & \textbf{0.667} & \textbf{12.7\,mm} & 1.62\\\bottomrule
\end{tabular}
\caption{EXP-M5-002, 150 paired test episodes. D2 beats D, C1 and C2 (each Holm $p=0.0078$);
D2 vs random $+0.033$, $p=0.54$.}
\end{table}
D2 reached the lowest centre error with fewer distinct views than random: it chooses views whose
biases cancel rather than simply moving more.

\subsection{Self-calibration from the collected views (EXP-M5-003)}
Correcting the grasp by the online bias estimate raised D2 from 0.640 to 0.767 ($+0.127$
[0.047, 0.207], $p=0.0034$). The fixed-camera control showed exactly zero effect, as predicted:
with one view there is nothing to identify. The full system beats fixed by $+0.253$ (0.513
$\rightarrow$ 0.767). Self-calibrated D2 and self-calibrated random are not distinguishable
($+0.020$, $p=0.75$).

\subsection{Closed-loop pick-and-place and recovery (EXP-M6-002)}
With a tray placement and up to two attempts, success is 0.50 (fixed) and 0.64 (D2). Retries
recovered few failures (14\,\% retry success; H6 not supported): a re-plan from the same belief
does not remove calibration bias.

\subsection{Final study and ablation (EXP-M8-001)}\label{sec:m8}
Eight arms ran on 300 new test seeds (299 paired; one seed had no confirmed track in any arm).
B0 and B1 share B2's views and target and differ only in the pose the grasp is planned from.
\begin{table}[h]\centering\small
\begin{tabular}{lc}\toprule
Arm & Success\\\midrule
B0 single-frame pose, fixed views & 0.495\\
B1 Kalman, fixed views & 0.488\\
B2 belief, fixed views & 0.482\\
B3 entropy views & 0.515\\
random views & 0.619\\
D2 calibrated task risk & 0.656\\
\textbf{D2 + self-calibration (full)} & \textbf{0.773}\\\bottomrule
\end{tabular}
\caption{EXP-M8-001. Primary family (Holm): full vs B1 $+0.284$ [0.217, 0.351]; D2 vs B3
$+0.140$ [0.080, 0.201]; full vs D2 $+0.117$ [0.070, 0.167] (all Holm $p<10^{-4}$); B2 vs B1
with fixed views $-0.007$, $p=0.85$ (not supported).}
\end{table}
Three of four primary hypotheses are confirmed. The fourth is a clean null: the belief's rotation
advantage (EXP-M4-001, Section~3) does not become grasp success while a calibration bias limits the grasp.
Estimator choice is worth at most 1.3 points; view choice is worth 17 and bias estimation 12.
D2 and random are again not distinguishable without self-calibration ($+0.037$, $p=0.31$).

\emph{Integrity note.} The self-calibration arms ran as a separate cluster array. Replay is
bit-identical on one machine (15/15 same-node pairs) but not across the cluster's heterogeneous
CPUs (75/284 cross-node pairs differ slightly). The fixed-camera control, whose correction is
numerically zero, measures this noise at 3 vs 8 flipped outcomes ($p=0.23$, no direction).
Comparisons within one array (P2, P4) are unaffected; P1 and P3 exceed it by an order of
magnitude.

\subsection{Uncertainty calibration (EXP-M8-002)}
We scale B2's reported position covariance by one scalar per noise preset, fitted on the
calibration split, and measure coverage on new test episodes (${\approx}2{,}670$ track-frames per
preset). Coverage is 0.54/0.80/0.92 (nominal) and 0.52/0.81/0.93 (harsh) at 50/80/95\,\%: all
six within $\pm5$\,pp, so H2 holds. We had predicted failure from a small val probe and report the
miss. The nominal 95\,\% interval [0.894, 0.934] excludes 0.95, so the tail is slightly
under-covered. Raw covariances are ${\sim}700\times$ too wide under nominal noise because they carry
a calibration allowance that the nominal preset lacks; a single global scale fails (harsh
0.33/0.59/0.83). The Kalman tracker fails the same test (nominal 0.69/0.82/0.87).

\subsection{Systems}
The C++ free-space likelihood matches NumPy to $2.7\times10^{-13}$ and makes a belief update
3.4$\times$ faster (164 $\rightarrow$ 48\,ms). On a V100 the pose proposer runs at 0.48\,ms per
call in ONNX Runtime (batch 1); TensorRT~10.16 does not support SM\,70.
On an A100, TensorRT fp16 takes 0.19\,ms at batch 1 (4.9$\times$ PyTorch fp32) and reaches
161k samples/s at batch 256, with top-1 error unchanged (4.921 vs 4.917\,mm). Perception remains ${\sim}$237\,ms per frame on
CPU, so the system is not real-time.

\section{Limitations}\label{sec:limits}
Simulation only; no hardware validation is claimed. One policy regime and one calibration-error
model; objects are primitives and mugs with known models. D2 and self-calibration were designed
after observing EXP-M5-001, though each was confirmed on new seeds. FoundationPose was not run.
The B0/B1 baselines in EXP-M8-001 share B2's target choice, isolating the pose estimator rather
than forming independent pipelines. Results replay exactly per machine but not across
heterogeneous cluster CPUs (Section~\ref{sec:m8}). Coverage calibration is per noise preset, which
assumes the deployment knows its sensor configuration.

\section{Conclusion}
A task-aware view objective helps only when its loss model contains the error that actually
limits the task. Here that error was view-dependent calibration bias, not pose ambiguity, and the
benefit comes mostly from view diversity plus estimating the bias it reveals.

\paragraph{Reproducibility.} Code, registry and reports:
\url{https://github.com/parammadan/belief-3d}.

\bibliographystyle{plainnat}
\bibliography{refs}
\end{document}
